\documentclass{iLRNConfProc}

\title{Contribution Title: Subtitle}

\author{
\begin{minipage}{\textwidth}
\centering
First Author\textsuperscript{1 [Add ORCID]} and Second Author\textsuperscript{2 [Add ORCID]} \\
\textsuperscript{1}Institution Name 1, City, Country \\
\textsuperscript{2}Institution Name 2, City, Country \\
corresponding\_author\_email@institution.com
\end{minipage}
}

\date{}

\begin{document}
\maketitle
\thispagestyle{empty}

\iLRNabstract{The abstract is a snapshot of your Short Paper / Extended Abstracts / Doctoral colloquium paper / \textit{iLEAD / iLRNFuser} contribution submission and should provide a summary that is between 150--250 words. Please use the Times New Roman 10pt font and justify the indents like this example. Accepted papers will appear in a proceedings book published by \textit{iLRN} with ISBN and individual DOIs.}

\iLRNkeywords{First Keyword, Second Keyword, Third Keyword.}

\section{First Section: Introduction}
Please use this formatting template to submit your proposal. It shows the fonts and styles required as part of the \textit{iLRN} guidelines for the Academic Stream (Work-in-progress and Doctoral Colloquium only) and Practitioner Stream conference proceedings.

\subsection{A Subsection Sample}
Note that the first paragraph of a section or subsection is not indented. Likewise, the first paragraph that follows a table or figure does not have an indent.

\indent Subsequent paragraphs, however, are indented as in this example. All printed material, including text, tables, and figures, must be within the margins of this template, which has a 1-inch (2.54 cm) margin all around. The body text should be justified, single-spaced, and Times New Roman, 10pt font.

Please ensure the accuracy of your submission. You must have all relevant permissions and clearances to present materials at \textit{iLRN} 2023. Organizers of the \textit{iLRN} conference accept no responsibility for statements made by authors either in written papers or in oral presentations. All authors retain copyright of their work.

\subsubsection*{Sample Heading (Third Level)}
Only two levels of headings should be numbered. Lower-level headings remain unnumbered; they are formatted as run-in headings.

\paragraph{Sample Heading (Fourth Level)}
The contribution should contain no more than four levels of headings. The following Table~\ref{tab:headings} gives a summary of all heading levels.

\begin{table}[h]
\centering
\caption{Table captions should be placed above the tables.}
\begin{tabular}{lll}
\specialrule{1.5pt}{0pt}{0pt}
\textbf{Heading level} & \textbf{Example} & \textbf{Font size and style} \\
\midrule
Title (centered) & \textbf{Main Title} & 14 \textit{point}, bold \\
1\textsuperscript{st}-level heading & \textbf{1 Title\_level\_1} & 12 \textit{point}, bold \\
2\textsuperscript{nd}-level heading & \textbf{1.1 Title\_level\_2} & 10 \textit{point}, bold \\
3\textsuperscript{rd}-level heading & \textbf{Title\_level\_3.} Text follows & 10 \textit{point}, bold \\
4\textsuperscript{th}-level heading & \textit{Title\_level\_4.} Text follows & 10 \textit{point}, italic \\
\specialrule{1.5pt}{0pt}{0pt}
\end{tabular}
\label{tab:headings}
\end{table}

You are encouraged to use tables and figures to represent your ideas. There is no limit on the number of tables or figures, but use good judgment. Submit high-resolution images (see Fig.~\ref{fig:example}).

\begin{figure}[htbp]
\centering
\includegraphics[width=0.5\textwidth]{example-image} % Replace with actual image path
\caption{A figure caption is always placed below the illustration. Short captions are centered, while long ones are justified.}
\label{fig:example}
\end{figure}

\section*{Acknowledgments (optional)}
Use this section to acknowledge additional contributors who do not satisfy authorship criteria, if applicable. Ensure you have obtained permission from those mentioned here. Please also recognize any grants that have supported your research.

\section*{Citations}
For reference citations, we prefer using square brackets and consecutive numbers. The following list provides a sample reference list with entries for journal articles \cite{ref1}, a conference paper \cite{ref2}, a book \cite{ref3}, proceedings without editors \cite{ref4}, and a URL \cite{ref5}. For more examples, see the Springer’s MathPhySci format.

\bibliographystyle{spmpsci}
%\bibliography{references}  %use a .bib file 

\begin{thebibliography}{5}  % or use thebibliography environment for manual entries

\bibitem{ref1}
Author, F.: Article title. \textit{Journal} \textbf{2}(5), 99--110 (2025)

\bibitem{ref2}
Author, F., Author, S.: Title of a proceedings paper. In: Editor, F., Editor, S. (eds.) \textit{CONFERENCE 2016}, ABCD, vol. 9999, pp. 1--13. Publisher, Location (2025)

\bibitem{ref3}
Author, F., Author, S., Author, T.: \textit{Book title}. 2nd ed. Publisher, Location (2025)

\bibitem{ref4}
Author, F.: Contribution title. In: \textit{9th International Proceedings on Proceedings}, pp. 1--2. Publisher, Location (2025)

\bibitem{ref5}
Webpage, \url{http://www.webapge.com}, last accessed 2025/04/01

\end{thebibliography}

\end{document}